\section*{第 \ref{ch:b3:gravity-geodesy} 章——重力与地球形状}

\begin{solution}{exo:b3:gravity-geodesy:1}
$b = a(1 - f) = 6378137\times0.99664719 = \qty{6356752.3}{m}$，且
$a - b = \qty{21384.7}{m}$，约 \qty{21.4}{km}。由克莱罗关系：
$f = (3J_2 + m)/2 = (0.00324789 + 0.00344979)/2 = \num{0.0033488}$，或
$1/298.61$，与 $1/298.26$ 相比小 \qty{0.12}{\percent}，此即
克莱罗关系所忽略的二阶项的量级。
\end{solution}

\begin{solution}{exo:b3:gravity-geodesy:2}
$\sigma = \qty{6.0e13}{kg}/\qty{5.0e11}{m^2} = \qty{120}{kg.m^{-2}}$，故
$h_w = 120/1000 = \qty{0.12}{m}$：相当于 \qty{12}{cm} 厚的水层。若将其
均匀分布于海洋，则海平面变化 $60/360 = \qty{0.17}{mm}$。
\end{solution}

\begin{solution}{exo:b3:gravity-geodesy:3}
在赤道处 $\gamma = \qty{9.780327}{m.s^{-2}}$。在 \ang{45} 纬度处，
$\sin^2\varphi = 0.5$，且 $\sin^22\varphi = 1$：$\gamma = 9.780327\times1.0026454
= \qty{9.806200}{m.s^{-2}}$。在两极处，$9.780327\times1.0053024 =
\qty{9.832186}{m.s^{-2}}$。两者之差为 \qty{0.051859}{m.s^{-2}}，即
\qty{5186}{\milli\gal}。
\end{solution}

\begin{solution}{exo:b3:gravity-geodesy:4}
$D = ET_e^3/(12(1 - \nu^2)) = 10^{11}T_e^3/11.25$：\qty{10}{km} 时为 \qty{8.89e21}{N.m}，
\qty{30}{km} 时为 \qty{2.40e23}{N.m}。取 $(\rho_m - \rho_w)g =
\qty{22269}{N.m^{-3}}$，$\alpha = (4D/22269)^{1/4}$：分别为 \qty{35.5}{km} 和
\qty{81.0}{km}。凹槽止于 $3\pi\alpha/4$，即 \qty{84}{km} 和 \qty{191}{km}；
\omterm{def:b3:gravity-geodesy:alpha}{前缘隆起}的峰顶位于 $\pi\alpha$，即 \qty{112}{km} 和 \qty{255}{km}。
\end{solution}

\begin{solution}{exo:b3:gravity-geodesy:5}
(a) $T = -\gamma N = -9.80\times50 = \qty{-490}{J.kg^{-1}}$。(b) 大地水准面相对于
椭球倾斜 $10/100\,000 = \num{1.0e-4}$ rad，铅垂线
与大地水准面正交，其相对于椭球法线同样倾斜相同角度：
$\num{1.0e-4}\times206\,265 = 20.6''$.
\end{solution}

\begin{solution}{exo:b3:gravity-geodesy:6}
$H = J_2/(C/Ma^2) = \num{1.08263e-3}/0.3308 = \num{0.0032728}$，且
$1/H = 305.55$。依题设 $J_2$ 不变；均匀地球应有
$C/(Ma^2) = 0.4$，故 $H = J_2/0.4 = \num{0.0027066}$，岁差将
变慢。同时测定 $J_2$ 与 $H$，
即可揭示质量向地心集中的程度。
\end{solution}

\begin{solution}{exo:b3:gravity-geodesy:7}
$r = \qty{6878.137}{km}$；$n = \sqrt{GM/r^3} = \qty{1.10678e-3}{rad.s^{-1}}$，
周期 $2\pi/n = \qty{94.6}{min}$。$\dot\Omega = -1.5\times\num{1.10678e-3}
\times\num{1.08263e-3}\times(6378.137/6878.137)^2\times\cos\ang{89} =
\qty{-2.70e-8}{rad.s^{-1}}$，即每天 \ang{-0.134}，向西。近极轨道
经过每一纬度，因此能测绘整个地球；其轨道亦
缓慢进动，使地面轨迹规律地填补空白。
\end{solution}

\begin{solution}{exo:b3:gravity-geodesy:8}
$e^{-2\pi\times500/100} = e^{-31.4} = \num{2.3e-14}$：\qty{2e-13}{\milli\gal}，
可忽略。$e^{-2\pi\times500/1000} = e^{-3.14} = 0.0432$：
\qty{0.43}{\milli\gal}。
\end{solution}

\begin{solution}{exo:b3:gravity-geodesy:9}
$3\pi\alpha/4 = \qty{120}{km}$ 给出 $\alpha = \qty{50.9}{km}$。
$D = (\rho_m - \rho_w)g\alpha^4/4 = 22269\times(50\,930)^4/4 =
\qty{3.75e22}{N.m}$，且 $T_e = (11.25D/E)^{1/3} = \qty{16.2}{km}$。板块
在距海沟 $\pi\alpha/2 = \qty{80}{km}$ 处穿过其未挠曲水平面。
\end{solution}

\begin{solution}{exo:b3:gravity-geodesy:10}
(a) 令 $u = x/\alpha$，$\mathrm d^4/\mathrm dx^4 = \alpha^{-4}\,\mathrm
d^4/\mathrm du^4$。由 $\cos u + \sin u = \sqrt2\cos(u - \pi/4)$，
函数为 $\sqrt2w_0\,\mathrm{Re}\,e^{(-1 + i)u - i\pi/4}$，且
$(-1 + i)^4 = -4$：故 $w'''' = -4w/\alpha^4$。于是 $Dw'''' + (\rho_m -
\rho_i)gw = (-4D/\alpha^4 + (\rho_m - \rho_i)g)w = 0$，因为 $\alpha^4 =
4D/((\rho_m - \rho_i)g)$。(b) $w'' = -(2w_0/\alpha^2)e^{-u}(\cos u - \sin u)$，故
$M = -Dw'' = (2Dw_0/\alpha^2)e^{-u}(\cos u - \sin u)$：绝对值在
$u = 0$ 处最大，零点位于 $\tan u = 1$，即 $x = \pi\alpha/4$。(c) 在载荷之下，
板块呈上凹，上表面受压；超过 $\pi\alpha/4$ 后
则呈上凸，上表面受拉。此处 $M' \propto
-e^{-u}\cos u$ 在 $u = \pi/2$ 处为零：拉伸在
$x = \pi\alpha/2$ 处最大，位于凹槽朝\omterm{def:b3:gravity-geodesy:alpha}{前缘隆起}一侧的外坡上，此处
$|M|$ 为其载荷下数值的 $e^{-\pi/2} = \qty{21}{\percent}$。正是
在此处，板块顶部的正断层与裂隙将会张开。
\end{solution}

\begin{solution}{exo:b3:gravity-geodesy:11}
$\sqrt{1 + \eta} = 2.5(1 - 1.5\times0.3308) = 1.2595$，故 $1 + \eta = 1.5863$
且 $\tfrac52m/f = 2.5863$：$f_h = 2.5\times0.00344979/2.5863 = \num{0.0033346}$，
$1/f_h = 299.9$。差值 $f - f_h = \num{1.82e-5}$；乘以 $a$，赤道
高出静力平衡形状 \qty{116}{m}。地球并非完全的
平衡流体：其内部承受应力，源自岩石圈的强度，
或地幔的流动与密度差异，这些应力维持
其形状略偏离平衡态。
\end{solution}

\begin{solution}{exo:b3:gravity-geodesy:12}
升交点须每秒推进 $2\pi/(365.242\times86\,400) =
\qty{1.991e-7}{rad.s^{-1}}$。当 $r = \qty{7078.137}{km}$ 时，$n =
\qty{1.0602e-3}{rad.s^{-1}}$，$\tfrac32nJ_2(a/r)^2 = \num{1.398e-6}$，故
$\cos i = -\num{1.991e-7}/\num{1.398e-6} = -0.1424$，$i = \ang{98.2}$：一个
略呈逆行的太阳同步轨道。
\end{solution}

\begin{solution}{pb:b3:gravity-geodesy:1}
\textbf{1.} $m = (\num{7.292115e-5})^2\times6378137^2\times6356752.3/
\num{3.986005e14} = \num{0.00344979}$.

\textbf{2.} $f = (3J_2 + m)/2 = \num{0.0033488}$, $1/f = 298.61$.

\textbf{3.} $\sqrt{1 + \eta} = 2.5(1 - 1.5\times0.3308) = 1.2595$,
$\tfrac52m/f_h = 2 + 0.5863$，故 $f_h = \num{0.0033346}$，且 $1/f_h = 299.88$。

\textbf{4.} 克莱罗值 $1/298.61$ 与
观测值 $1/298.26$ 相差不超过 \qty{0.12}{\percent}；静力学值 $1/299.88$ 则
偏小 $\Delta f = \num{1.82e-5}$，故 $a\,\Delta f = \qty{116}{m}$：赤道
比静力学形状高出 \qty{116}{m}。

\textbf{5.} $H = \num{1.08263e-3}/0.3308 = \num{0.0032728}$ ($1/305.6$).

\textbf{6.} $\tfrac12\times\qty{80}{km}\times\qty{6}{km} = \qty{240}{km^2} =
\qty{2.40e8}{m^2}$.

\textbf{7.} $V_0 = (2800 - 1030)\times9.81\times\num{2.40e8} =
\qty{4.17e12}{N.m^{-1}}$.

\textbf{8.} 艾里模型：$(\rho_m - \rho_v)r = (\rho_v - \rho_w)h$，故 $r =
6\times1770/500 = \qty{21.2}{km}$，位于山脊之下。

\textbf{9.} 海底远远超出山脊基部发生挠曲：
载荷并非逐柱补偿，而是由板块在
宽四倍的区域上弯曲来补偿。板块具有强度：这就是区域补偿。

\textbf{10.} $\alpha = 165/(3\pi/4) = \qty{70.0}{km}$.

\textbf{11.} $\alpha = 220/\pi = \qty{70.0}{km}$：结果相同，正如单一弹性
板块所要求的。

\textbf{12.} $D = (\rho_m - \rho_w)g\alpha^4/4 = 2270\times9.81\times(\num{7.00e4})^4/4
= \qty{1.34e23}{N.m}$.

\textbf{13.} $T_e = (12(1 - \nu^2)D/E)^{1/3} = (11.25\times\num{1.34e23}/
\num{1e11})^{1/3} = \qty{24.7}{km}$.

\textbf{14.} $w_0 = V_0\alpha^3/(8D) = \num{4.17e12}\times(\num{7.00e4})^3/
(8\times\num{1.34e23}) = \qty{1334}{m}$.

\textbf{15.} $w_0e^{-\pi} = 1334\times0.0432 = \qty{58}{m}$.

\textbf{16.} $2\times1334\times\num{7.00e4}\times2270\times9.81 =
\qty{4.16e12}{N.m^{-1}}$，与 $V_0$ 相等，误差在 $D$ 的舍入范围之内。

\textbf{17.} \qty{24.7}{km} 落在夏威夷海岭的数值范围之内。
板块在轴线下方下沉 \qty{1.3}{km}，而不是
形成 \qty{21}{km} 的山根：同样的浮力分布于沟槽的 \qty{330}{km} 范围内。

\textbf{18.} $\lambda = 2\pi\times\qty{70.0}{km} = \qty{440}{km}$,
$k = 1/\alpha = \qty{1.43e-5}{m^{-1}}$.

\textbf{19.} $\sigma = (2900 - 1030)\times500 = \qty{9.35e5}{kg.m^{-2}}$，故
$2\pi G\sigma = \qty{3.92e-4}{m.s^{-2}} = \qty{39.2}{\milli\gal}$，在海底处；
乘以 $e^{-k\times\qty{5}{km}} = 0.931$，在海平面处为 \qty{36.5}{\milli\gal}
（板块被压低之处为负值）。

\textbf{20.} $\sigma = (3300 - 2900)\times500 = \qty{2.0e5}{kg.m^{-2}}$：
\qty{8.39}{\milli\gal}，乘以 $e^{-k\times\qty{12}{km}} = 0.843$，
得 \qty{7.07}{\milli\gal}。总和为 \qty{43.6}{\milli\gal}。

\textbf{21.} $e^{-k\times\qty{400}{km}} = e^{-5.71} = 0.0033$：\qty{0.14}{\milli\gal}。

\textbf{22.} $e^{-2\pi\times400/80} = e^{-31.4} = \num{2.3e-14}$：卫星
看不到海岭本身，只能看到宽广的挠曲，而即便如此，
挠曲信号也被削弱了三百倍；海岭自身的异常须在海上测量。

\textbf{23.} $n \approx 2\pi\times6371/440 = 91$.

\textbf{24.} 由第 10 至 13 题可知，奥鲁海岭下方板块的
\omterm{def:b3:gravity-geodesy:flexure}{有效弹性厚度}为
\textbf{\qty[text-series-to-math=true]{24.7}{km}}.
\end{solution}
